\section{实验环境与数据}

\subsection{官方数据来源}

实验使用 OpenCV 官方仓库 \texttt{opencv/opencv} 的 \texttt{4.x/samples/data} 目录中的
\texttt{aloeL.jpg} 和 \texttt{aloeR.jpg}。选择该图像对的原因是：它属于 OpenCV 官方样例，左右文件
关系明确，官方 Python 立体匹配示例也针对 aloe 图像给出了可复现的 SGBM 处理思路，
适合在不引入额外数据集的情况下完成课程实验。

两幅 JPG 原始尺寸均为 $1282\times1110$，均为三通道图像。程序默认对左右图像同步执行
一次 \texttt{pyrDown}，实际匹配尺寸为 $641\times555$。同一官方目录中的 \texttt{aloeGT.png} 是
可选的视差参考图；本实验没有把它当作标定参数、深度尺度或误差真值。

\begin{table}[H]
  \centering
  \caption{实验数据与软件环境}
  \label{tab:environment}
  \begin{tabular}{ll}
    \toprule
    项目 & 实测版本或说明 \\
    \midrule
    数据目录 & OpenCV 官方 \texttt{4.x/samples/data} \\
    输入文件 & \texttt{aloeL.jpg}、\texttt{aloeR.jpg} \\
    原始尺寸 & $1282\times1110$，三通道 \\
    Python & 3.13.5 \\
    OpenCV runtime & 5.0.0 \\
    NumPy & 2.5.3 \\
    \bottomrule
  \end{tabular}
\end{table}

版本信息、输入路径和运行参数来自程序生成的
\codepath{task2/results/params.json}。输入文件的 SHA-256 与下载 URL 已记录在
\codepath{task2/results/run_evidence.md}，便于在同一版本环境中复核，而不依赖报告中
手工复制的图像统计。

\subsection{实际运行命令}

本次结果由以下默认命令生成，未使用 \texttt{--no-pyrdown}，也未额外调整匹配参数：

\begin{lstlisting}[language=bash]
task2/.venv/bin/python task2/src/stereo_depth.py \\
  --left task2/data/aloeL.jpg \\
  --right task2/data/aloeR.jpg \\
  --out-dir task2/results
\end{lstlisting}

程序先检查左右图像均可读取且尺寸一致，再进行灰度转换、同步降采样、StereoSGBM
匹配、视差筛选和相对深度计算。实验不使用其他图像对的标定文件，所有结果均来自
上述命令对应的根目录输出。
